\uuid{VQzy}
\titre{Théorème de Bézout}
\auteur{bodin ; croizat ; sacre}
\organisation{exo7}
\licence{CC BY-NC-SA 4.0}
\datecreate{2022-03-08}
\niveau{L1}
\module{Algèbre}
\chapitre{Arithmétique dans Z}
\sousChapitre{Nombres premiers, nombres premiers entre eux}
\difficulte{3}
\theme{Théorème de Bézout}
\exo7qcmid{hilisit.arithmetique.024}

\contenu{
\begin{question}
\qclassification{103.04}
    Quelles sont les affirmations vraies concernant l'algorithme d'Euclide ?
    \begin{answers}
        \bad{Il se peut que le processus n'aboutisse pas à cause d'un nombre infini de divisions à effectuer.}
        \bad{Il se peut que le processus ne fournisse pas le pgcd correct.}
        \good{Le pgcd est le dernier reste non nul.}
        \good{L'algorithme étendu permet en plus de calculer des coefficients de Bézout.}
    \end{answers}
    \begin{explanations}
        L'algorithme d'Euclide fournit un résultat \emph{correct} en un nombre \emph{fini} d'étapes. La remontée de l'algorithme d'Euclide permet de calculer des coefficients de Bézout.
    \end{explanations}
\end{question}
}
